\question[4]
Ein Unternehmen entwickelt ein neues Verschlüsselungsverfahren. Die Entwickler
sind der Meinung, dass das Verfahren sicher ist, weil der verwendete Algorithmus
nicht veröffentlicht wird.

Erläutere das Kerckhoffs’sche Prinzip und bewerte die Aussage der Entwickler.
Gehe dabei darauf ein, wovon die Sicherheit eines modernen
Verschlüsselungsverfahrens tatsächlich abhängen sollte.

\begin{solutionbox}{4cm}
Nach dem Kerckhoffs’schen Prinzip darf die Sicherheit eines
Verschlüsselungsverfahrens nicht von der Geheimhaltung des Algorithmus abhängen.
Es wird angenommen, dass Angreifer den Algorithmus kennen, 
geheim bleiben soll lediglich der verwendete Schlüssel.
Ein sicheres Verfahren muss auch dann sicher bleiben, wenn sein Aufbau 
öffentlich bekannt ist. Die Aussage der Entwickler ist daher nicht überzeugend
und widerspricht dem Kerckhoffs’schen Prinzip.
\end{solutionbox}